\clearpage\fancyhead[L]{\small AUSZUG AUS EINEM BRIEF.}\hypertarget{native-323-lower}{}\label{native-323-lower}

\begin{center}\textbf{AUSZUG AUS EINEM BRIEF VON PROFESSOR GAUSS.}\par\vspace{4pt}\hrule\vspace{4pt}\fontsize{10}{12}\selectfont Monatliche Correspondenz. Band XVIII. S. 289. 290. 1810 September.\par\vspace{4pt}\hrule\vspace{4pt}\end{center}
\begin{flushright}Göttingen, den 17. August 1810.\end{flushright}

\par\indent ---\enspace---\enspace--- Sie werden bereits wissen, dass unser König jetzt den Weiterbau unserer neuen Sternwarte angeordnet und dafür mit wahrhaft königlicher Freigebigkeit 200000 Franken bewilligt hat
\SourceLeaf{324}Wenn Sie diesen Herbst noch zu uns kommen, werden Sie hoffentlich sehen, dass der Bau bereits wieder aufgenommen wurde.

Am 8. August war ich in Münden, unter anderem um herauszufinden, wie sich die Fahrt auf den Gang unseres Chronometers auswirkt. Das Wetter war jedoch sehr ungünstig, und vormittags zwischen 10 und 11 Uhr konnte ich nur einige ungenaue Sonnenhöhen messen. Mittags war die Sonne völlig unsichtbar, und erst nach 1 Uhr erhielt ich einige zuverlässige Höhen. Aus diesen und einigen ebenso guten Messungen gegen 6 Uhr bestimmte ich die Polhöhe mit $51^\circ\ 25'\ 22''$ und den Längenunterschied gegenüber \emph{Göttingen} zu $1^{\text{m}}9^{\text{s}}3$ westlich, in Zeit gemessen. Der Chronometer hat sich ausgezeichnet bewährt. Sein täglicher Gang vom 7\,..\,9. August betrug $0^{\text{s}}1$ und war damit derselbe wie zuvor. Da meine Sonnenhöhen so weit vom Mittag entfernt gemessen wurden, möchte ich die Genauigkeit meiner Bestimmung erst beurteilen, wenn ich Gelegenheit habe, sie zu wiederholen. Ich muss jedoch bemerken, dass meine neun Sonnenhöhen gut miteinander übereinstimmen. Mein Beobachtungsplatz war der sogenannte \emph{Freitagswerder} am Zusammenfluss von \emph{Fulda} und \emph{Werra} unmittelbar nördlich vor der Stadt. Professor \textsc{Seyffer} bestimmte 1794 an einem vermutlich etwas \emph{südlicher} gelegenen Ort die Polhöhe mit $51^\circ\ 26'\ 52''$, allerdings nur aus einer einzigen Höhe.
